\subsection{Math Exploration \& Vibe Coding Agents}
\label{sec:app-math}
Mathematics and code have traditionally served as two of the most widely used domains for evaluating reasoning in artificial intelligence, as both require structured symbolic manipulation and precise multi-step deduction. Traditional benchmark-driven evaluation in these domains is showing clear limitations. Widely used math datasets such as GSM8K \cite{cobbe2021training}, MATH \cite{hendrycks2021measuring}, and AIME \cite{MAA_AIME} are increasingly saturated, which makes it difficult to distinguish among modern high-performing models. The problems in these datasets often rely on a small set of recurring techniques and do not require the sustained and exploratory reasoning needed to assess more advanced mathematical capabilities. Even recent evaluations such as FrontierMath \cite{glazer2024frontiermath} continue to emphasize final-answer accuracy, which offers only a partial view of an agent’s reasoning process and its ability to adjust strategies during problem solving.

Under the agentic reasoning paradigm, however, both areas are undergoing a substantial shift from static problem solving to dynamic processes that emphasize exploration, adaptation, and collaboration. In mathematics, recent systems \citep{novikov2025alphaevolve,luong2025towards,trinh2024solving,romera2024mathematical} demonstrate that agents can engage in competition-level reasoning, building on the success of LLMs in coding tasks. Work in foundational mathematics \citep{swirszcz2025advancing,georgiev2025mathematical} further shows that agents can search for new problems, propose conjectures, construct auxiliary lemmas, and explore deeper structures in mathematical concepts. These developments position mathematics not merely as an evaluation benchmark but as a domain of active \emph{mathematical exploration}.

Large Language Models have also reshaped coding through the emerging workflow known as agentic coding and \emph{vibe coding} \citep{karpathy2025vibecoding,willison2025vibecoding}. In this paradigm, the model acts as an interactive collaborator that engages in multi-turn natural-language dialogue. Users iteratively design and refine programs while the agent maintains context, adapts to evolving requirements, and continuously self-corrects. Modern tools such as Copilot\footnote{\url{https://github.com/features/copilot}} and Cursor\footnote{\url{https://cursor.com}} have further popularized this collaborative workflow, making interactive programming a common practice in real-world software development.

In this section, we organize our discussion according to the three-layer framework introduced earlier. The foundational layer (Section~\ref{sec:app-math-planning}) concerns the core reasoning and execution skills: mathematical agents perform symbolic manipulations and step-by-step derivations across arithmetic, algebra, geometry, and calculus, while code agents carry out syntax-aware generation, implement functions, and verify correctness through interpreter or compiler feedback. The self-evolving layer (Section~\ref{sec:app-math-evolving}) introduces mechanisms for reflection and adaptation. Mathematical agents learn from intermediate reasoning traces to correct missteps or explore alternative solution paths, and code agents iteratively debug, refine, and optimize implementations based on runtime feedback or test results. The collective layer (Section~\ref{sec:app-math-collaboration}) focuses on collaboration, where agents exchange intermediate results, share reusable modules, and jointly develop complex proofs or codebases. Taken together, these layers reveal how mathematics and coding are becoming domains in which agentic reasoning enables increasingly creative and adaptive problem solving.


\subsubsection{Foundational agentic reasoning}
\label{sec:app-math-planning}


\paragraph{Planning.}

Explicit planning is widely recognized as a core mechanism for enhancing the structured 
reasoning capabilities of LLMs. In the domain of mathematical discovery, several systems exhibit structures that can be interpreted as forms of planning. In representation theory and knot theory, the system of \citet{davies2021advancing} guides human mathematicians by proposing intermediate objects and promising avenues of exploration, which function as high-level suggestions for organizing problem-solving workflows. In geometric reasoning, \citet{trinh2024solving} solves Olympiad-level geometry problems by decomposing them into sequential stages of construction, lemma generation, and verification, yielding a structured multi-step process that resembles a planned reasoning trajectory. Program-search approaches \citep{romera2024mathematical} iteratively refine candidate programs and mathematical structures, a procedure that naturally forms a coarse-to-fine exploration path. Large-scale exploration frameworks \citep{georgiev2025mathematical,swirszcz2025advancing} also operate through cycles of proposing, testing, and modifying conjectures or geometric objects, which collectively create a procedural structure aligned with planning. Efforts toward more robust mathematical reasoning \citep{luong2025towards} similarly rely on stepwise reasoning patterns, further reinforcing the presence of implicit planning dynamics. of implicit planning dynamics across mathematical agents.

In code agents, planning has likewise emerged as an essential component for organizing 
multi-step reasoning and enabling more structured decision-making. Early systems such as 
CodeChain \citep{le2023codechain} and CodeAct \citep{wang2024executable} introduce explicit planning 
or action spaces to support modular code construction, while KareCoder \citep{huang2024knowledge} 
integrate external knowledge sources or domain-specific information 
into the planning process. Subsequent works explore more structured planning organizations, including 
multi-stage control flows \citep{bairi2024codeplan,multi-stage}, tree-shaped planning structures 
\citep{li2024codetree,ni2024tree}, and adaptive refinement mechanisms \citep{aggarwal2025dars}. 
Planning has also been linked to improved exploration breadth: GIF-MCTS \citep{dainese2024generating} 
incorporates Monte Carlo Tree Search to explore multiple code-generation trajectories.
Recent extensions demonstrate applicability in specialized domains such as hardware design, 
where VerilogCoder \citep{ho2025verilogcoder} employs graph-structured planning and waveform-based 
verification. To address environments where state serialization is difficult, Guided Search 
\citep{zainullina2025guidedsearchstrategiesnonserializable} introduces lookahead and trajectory 
selection strategies for evaluating candidate actions without full environment access.





\paragraph{Tool-Use.}
Integrating external computational tools with LLMs has become a central mechanism for extending the reasoning and generation capabilities of single-agent systems. A defining characteristic of many mathematical reasoning systems is their integration with external computational tools. Formal theorem-proving agents such as \citet{thakur2024context} operate directly within the Lean proof assistant, selecting tactics and interacting with the underlying prover through in-context guidance. Position papers on formal mathematical reasoning \citep{yang2024formal} emphasize that progress in mathematical AI will depend on systems that can call theorem provers, satisfiability solvers, and computer algebra systems as part of a broader reasoning loop. Program-search frameworks for discovery \citep{romera2024mathematical} rely on executing generated programs and employing symbolic routines for verification. Generative modelling approaches \citep{ellenberg2025generative} make use of computational number-theoretic tools to check and filter generated candidates. Geometry-focused systems \citep{trinh2024solving,hubert2024ai} integrate automated geometric solvers and checkers to validate constructions and derived relations. Across these systems, external computational resources play a central role in enabling correct and scalable mathematical reasoning.

In code agents, external tools have similarly become crucial for extending the capabilities of 
LLM-based agents beyond pure text generation. Early work such as Toolformer \citep{schick2023toolformer} 
and ToolCoder \citep{zhang2023toolcoder} explored how models can learn to invoke APIs or search tools 
to obtain missing information during generation. Subsequent systems integrate increasingly rich 
toolchains: ToolGen \citep{wang2024toolgen} leverages automatic completion tools to resolve undefined 
dependencies, while CodeAgent \citep{zhang2024codeagent} incorporates multiple programming utilities 
 including search, documentation reading, symbol navigation, and code execution to support more 
realistic software workflows. Several methods focus on improving tool-feedback loops, such as ROCODE 
\citep{jiang2024rocode}, which combines real-time error detection with adaptive backtracking, and 
CodeTool \citep{lu2025codetool}, which introduces process-level supervision to improve the reliability 
of tool invocation. Collectively, these systems show that tool integration provides essential external 
signals, via search results, documentation, static analysis, or execution feedback, that extend the 
reasoning and generation capabilities of single-agent LLMs.

\paragraph{Search and Retrieval.}

Search and retrieval has emerged as a complementary mechanism that enriches model contexts through external information sources.
Search is a recurring mechanism in mathematical discovery. Program-search based systems \citep{romera2024mathematical} treat mathematical discovery as navigating a program space in which candidate programs encode conjectures or structural hypotheses, with iterative filtering based on symbolic or numerical checks. Generative modelling approaches \citep{ellenberg2025generative} explore families of mathematical objects by sampling from flexible distributions that capture structural regularities. Geometric systems such as \citet{trinh2024solving} and \citet{swirszcz2025advancing} search over constructions, configurations, and high-dimensional polytopes, guided by learned heuristics or structural constraints. Large-scale discovery frameworks \citep{georgiev2025mathematical} operate through repeated propose--test--refine cycles across conjectures, supporting wide exploration over mathematical landscapes. All these systems rely on systematic search procedures that structure the exploration of mathematical ideas.

In code generation, repository-level retrieval systems 
such as RepoHyper \citep{phan2024repohyper} locate reusable code segments from large-scale code bases 
to provide more informative contexts for generation. CodeNav \citep{gupta2024codenav} dynamically 
indexes real repositories during generation, retrieving relevant functions and adjusting based on 
execution feedback. AUTOPATCH \citep{acharya2025optimizing} applies retrieval to performance 
optimization, combining historical code examples with control flow graph analysis for 
context-aware improvements. Structure-aware retrieval has also been explored: knowledge-graph-based 
repository representations \citep{athale2025knowledge} improve retrieval quality by capturing 
symbolic and relational structure, while cAST \citep{zhang2025cast} introduces AST-based chunking to 
enhance syntactic coherence and retrieval granularity. These retrieval methods demonstrate how 
external knowledge sources can augment single-agent LLMs by providing high-quality, structured 
contexts that guide both understanding and generation.

\subsubsection{Self-evolving agentic reasoning}
\label{sec:app-math-evolving}

\paragraph{Agentic Feedback and Reflection.}
Across mathematical and code reasoning tasks, feedback operates as an external signal that 
highlights discrepancies, confirms correct inferences, and directs the agent toward more reliable 
subsequent computations. Feedback mechanisms appear prominently across mathematical discovery systems. In program-search 
based discovery \citep{romera2024mathematical}, executing candidate programs and evaluating their 
outputs against constraints yields counterexamples or confirmations, enabling iterative refinement of 
conjectures. In geometry, automated checkers validate constructions and derived relationships 
\citep{trinh2024solving,hubert2024ai}, providing correctness signals that guide subsequent revisions. 
Interactive evaluation frameworks \citep{collins2024evaluating} show that human clarifications and 
follow-up prompts expose reasoning errors and improve model responses. Position work on formal 
reasoning \citep{yang2024formal} highlights verification, proof checking, and model checking as 
essential sources of structured feedback. In several systems involving multiple candidate hypotheses 
\citep{romera2024mathematical,ellenberg2025generative,georgiev2025mathematical}, 
the use of verification signals to retain promising candidates functions analogously 
to a fitness-based evaluation step, since these signals determine which hypotheses 
survive and which are discarded, thereby shaping the direction of subsequent exploration 
without introducing an explicit learning signal.

For code agents, feedback and reflection are central to improving reliability over multi-step 
reasoning. Fault-aware editing methods such as Self-Edit \citep{zhang2023self} incorporate 
execution-based signals to refine erroneous code, while Self-Repair 
\citep{olausson2024selfrepairsilverbulletcode} integrates code and feedback models to diagnose test 
failures and propose targeted corrections. More structured systems like LeDeX \citep{jiang2024ledex} 
combine stepwise annotation, execution-driven verification, and automated repair into a closed-loop 
pipeline in which feedback continually informs the next revision. Reflection also functions as a 
form of memory: iterative self-improvement frameworks such as Self-Refine \citep{madaan2023self}, 
Self-Iteration \citep{chang2023self}, and Self-Debug \citep{chen2023teaching} reuse earlier drafts, 
analyses, and explanations to guide subsequent revisions, while artifact-level mechanisms such as 
CodeChain \citep{le2023codechain} and LeDeX \citep{jiang2024ledex} retain reusable components, 
corrected snippets, and execution traces as persistent representations. Together, these approaches 
demonstrate how feedback—whether symbolic, execution-based, or self-generated—interacts with 
iterative memory to support structured refinement and long-horizon improvement in code-oriented 
agentic systems.


\paragraph{Memory.}
Memory provides agents with a mechanism for retaining and leveraging information from earlier 
reasoning steps, allowing them to maintain consistency, improve intermediate states, and improve 
their performance over extended problem-solving horizons. While few systems introduce an explicit memory module, many mathematical agents rely on forms of persistent state that can be viewed as implicit memory. Interactive evaluation frameworks \citep{collins2024evaluating} maintain conversational and problem-state context across multiple turns, allowing models to build upon earlier partial derivations. Formal-theorem-proving agents \citep{thakur2024context} operate over evolving proof states in Lean, which accumulate tactics, subgoals, and intermediate lemmas, functioning as structured persistent information. Program-search and discovery systems \citep{romera2024mathematical,georgiev2025mathematical} retain conjecture histories, counterexamples, and successful constructions as part of their iterative refinement processes. Their role in preserving and reusing information across reasoning steps aligns with the broader notion of memory in agentic systems.



In code agents, memory increasingly takes the form of explicit structures that maintain coherence
over long-horizon generation. Several systems construct shared or structured workspaces: 
Self-Collaboration \citep{dong2024self} introduces a blackboard memory for storing task descriptions,
intermediate drafts, and revision records, enabling agents to coordinate through a common 
representation. Architectural approaches such as L2MAC \citep{holt2023l2mac} and Cogito 
\citep{li2025cogito} extend this idea by organizing context into dedicated registers, hierarchical 
memory units, or long-term knowledge stores, overcoming context-window limits and supporting 
multi-file or large-function reasoning. Across these designs, the underlying insight is consistent: effective code agents require 
persistent, structured, and often domain-aware memory that preserves intermediate reasoning and 
enables self-improvement across extended development trajectories.





\subsubsection{Collective multi-agent reasoning}
\label{sec:app-math-collaboration}
To address the growing complexity of tasks in mathematical discovery and code generation, recent 
systems increasingly rely on multi-agent or modular designs that decompose problems into 
cooperating specialized components. Mathematical discovery frameworks often organize reasoning into explicitly defined 
multi-agent or multi-component workflows that collaborate to explore and validate mathematical ideas. The polytope-generation system \citep{swirszcz2025advancing} uses multiple specialized components that generate, evaluate, and refine geometric objects, forming a genuine collaborative workflow. Large-scale exploration frameworks \citep{georgiev2025mathematical} often divide discovery into modules for proposing conjectures, identifying counterexamples, and refining statements, which, although implemented within a unified system, mirror multi-agent role specialization. Early work on AI-assisted mathematical research \citep{davies2021advancing} and Olympiad-level systems \citep{hubert2024ai} also involve human--AI collaboration, where human mathematicians interact with AI systems in a complementary manner. These developments indicate that mathematical discovery is an inherently collaborative process, and multi-agent architectures provide a natural vehicle for expressing such collaboration in agentic systems.

Multi-agent systems for code generation have progressed from simple role-based pipelines to adaptive, 
collaborative frameworks capable of handling long-horizon software development. Early approaches such 
as Self-Collaboration \citep{dong2024self} and AgentCoder \citep{huang2023agentcoder} decompose tasks 
into sequential roles, while hierarchical designs like PairCoder \citep{zhang2024pair} and FlowGen 
\citep{lin2024soen} introduce an architecture in which high-level agents handle planning and lower-level 
agents carry out concrete implementation. Flexible systems such as SoA \citep{ishibashi2024self} further 
adjust the number and specialization of agents in response to task complexity. Other frameworks, 
including MapCoder \citep{islam2024mapcoder}, AutoSafeCoder 
\citep{nunez2024autosafecoder}, and QualityFlow 
\citep{hu2025qualityflow}, rely on repeated cycles in which multiple agents 
generate, test, analyze, and repair code. Recent work explores self-evolving system structures, as in 
SEW \citep{liu2025sew}, which reorganizes collaboration pathways based on runtime feedback, and EvoMAC 
\citep{hu2024self}, which adjusts agent strategies through an iterative text-based update mechanism. 
Collaborative optimization methods such as Lingma SWE-GPT \citep{ma2024lingmaswegptopendevelopmentprocesscentric}, 
CodeCoR \citep{pan2025codecor}, SyncMind \citep{guo2025syncmind}, and CANDOR 
\citep{xu2025hallucinationconsensusmultiagentllms} explicitly improve cross-agent coordination. 
Together, these systems show a clear shift toward multi-agent code generators that rely not only on 
role decomposition, but also on reflection, distributed evaluation, adaptive restructuring, and 
team-level optimization, transforming code generation into an increasingly coordinated and resilient 
problem-solving process.

